\subsection{** APS-TSLCA-TSL-ACTIVATION-EQUATION **}\label{aps-tslca-tsl-activation-equation}

\subsection{** Three-Squared-Lattice Cognitive Architecture **}\label{three-squared-lattice-cognitive-architecture}

\subsection{** Symbiotic Universal Xessability Standards **}\label{symbiotic-universal-xessability-standards}

\subsection{** Aurphyx Primordial Standard **}\label{aurphyx-primordial-standard}

\subsection{** Aurphyx LLC **}\label{aurphyx-llc}

\subsection{** SAGES \textbar{} Proprietary \textbar{} Pro-Existence **}\label{sages-proprietary-pro-existence}

\subsection{** Accessibility = Xessability **}\label{accessibility-xessability}

\subsection{** Version 3.69 **}\label{version-3.69}

\section{Three‑Squared‑Lattice Activation Equation}\label{threesquaredlattice-activation-equation}

\emph{Rooted in the Triple Threshold Architecture and governed by the Harmonic Integrity Field}

\subsection{1. Lattice Structure}\label{lattice-structure}

The Three‑Squared‑Lattice is indexed as:

\[
\mathcal{L} = \{\, N_{i,j,k} \mid i,j,k \in \{1,2,3\} \,\}
\]

Each node \(N_{i,j,k}\) carries three local field values:

\begin{itemize}
\tightlist
\item
  \(C_{i,j,k}\) --- local coherence\\
\item
  \(R_{i,j,k}\) --- local resonance\\
\item
  \(A_{i,j,k}\) --- local alignment
\end{itemize}

Each node also inherits:

\begin{itemize}
\tightlist
\item
  its layer (Creation, Integration, Renewal)\\
\item
  its aXis (Coherence, Resonance, Alignment)\\
\item
  its mode (Local, Domain, Continuum)
\end{itemize}

This gives each node a \textbf{triple context} and a \textbf{triple field state}.

\begin{center}\rule{0.5\linewidth}{0.5pt}\end{center}

\subsection{2. Node‑Level HIF}\label{nodelevel-hif}

Each node computes its own Harmonic Integrity Field:

\[
\text{HIF}_{i,j,k} = \sqrt[3]{C_{i,j,k} \cdot R_{i,j,k} \cdot A_{i,j,k}} \cdot \Phi_{i,j,k}
\]

with the local coupling function:

\[
\Phi_{i,j,k} =
\begin{cases}
1 & \text{if } C_{i,j,k} \ge C_{\theta},\; R_{i,j,k} \ge R_{\theta},\; A_{i,j,k} \ge A_{\theta} \\
0 & \text{otherwise}
\end{cases}
\]

This is the \textbf{Triple Threshold Gate} at the node level.

A node is ``lit'' only when all three torches are lit.

\begin{center}\rule{0.5\linewidth}{0.5pt}\end{center}

\subsection{3. Neighborhood Coupling}\label{neighborhood-coupling}

Each node interacts with its \textbf{orthogonal neighbors}:

\[
\mathcal{N}(i,j,k) = \{\, N_{i\pm1,j,k},\; N_{i,j\pm1,k},\; N_{i,j,k\pm1} \,\}
\]

The neighborhood HIF is:

\[
\text{HIF}^{\text{nbr}}_{i,j,k} = \frac{1}{|\mathcal{N}|} \sum_{N \in \mathcal{N}(i,j,k)} \text{HIF}_N
\]

This captures local coherence circulation and resonance propagation.

\begin{center}\rule{0.5\linewidth}{0.5pt}\end{center}

\subsection{4. Layer‑Level HIF}\label{layerlevel-hif}

Each of the three layers computes:

\[
\text{HIF}^{\text{layer}}_{\ell} = \frac{1}{9} \sum_{(i,j,k) \in \ell} \text{HIF}_{i,j,k}
\]

where each layer contains 9 nodes.

This determines whether the \textbf{Creation}, \textbf{Integration}, or \textbf{Renewal} layer is active.

\begin{center}\rule{0.5\linewidth}{0.5pt}\end{center}

\subsection{5. Lattice‑Level HIF}\label{latticelevel-hif}

The global field value is:

\[
\text{HIF}^{\text{lattice}} = \frac{1}{27} \sum_{i,j,k} \text{HIF}_{i,j,k}
\]

This is the value used for:

\begin{itemize}
\tightlist
\item
  lattice‑scale creation\\
\item
  lattice‑scale integration\\
\item
  lattice‑scale renewal\\
\item
  governance legitimacy\\
\item
  cross‑domain compatibility
\end{itemize}

It is the \textbf{governing invariant} of the entire architecture.

\begin{center}\rule{0.5\linewidth}{0.5pt}\end{center}

\subsection{6. The Activation Equation}\label{the-activation-equation}

A node activates when \textbf{both} its local and neighborhood HIF exceed threshold:

\[
\Psi_{i,j,k} =
\begin{cases}
1 & \text{if } \text{HIF}_{i,j,k} \ge H_{\theta} \;\text{and}\; \text{HIF}^{\text{nbr}}_{i,j,k} \ge H_{\theta}^{\text{nbr}} \\
0 & \text{otherwise}
\end{cases}
\]

This is the \textbf{Three‑Squared‑Lattice Activation Equation}.

It encodes:

\begin{itemize}
\tightlist
\item
  local integrity\\
\item
  relational integrity\\
\item
  threshold integrity
\end{itemize}

A node cannot activate in isolation.\\
It must be coherent with its neighbors.

This is the computational mirror of the Hecate Triple Threshold Protocol.

\begin{center}\rule{0.5\linewidth}{0.5pt}\end{center}

\subsection{7. Propagation Dynamics}\label{propagation-dynamics}

Once activated, a node propagates its influence through:

\subsubsection{7.1 Coherence Propagation}\label{coherence-propagation}

\[
C_{i,j,k}(t+1) = f_C(C_{i,j,k}(t), \mathcal{N})
\]

\subsubsection{7.2 Resonance Propagation}\label{resonance-propagation}

\[
R_{i,j,k}(t+1) = f_R(R_{i,j,k}(t), \mathcal{N})
\]

\subsubsection{7.3 Alignment Propagation}\label{alignment-propagation}

\[
A_{i,j,k}(t+1) = f_A(A_{i,j,k}(t), \mathcal{N})
\]

The functions \(f_C, f_R, f_A\) are domain‑specific but must preserve invariants.

\begin{center}\rule{0.5\linewidth}{0.5pt}\end{center}

\subsection{8. Lattice State Transitions}\label{lattice-state-transitions}

The lattice enters:

\subsubsection{Creation Mode}\label{creation-mode}

\[
\text{HIF}^{\text{lattice}} \ge H_{\text{create}}
\]

\subsubsection{Integration Mode}\label{integration-mode}

\[
H_{\text{integrate}} \le \text{HIF}^{\text{lattice}} < H_{\text{create}}
\]

\subsubsection{Renewal Mode}\label{renewal-mode}

\[
\text{HIF}^{\text{lattice}} < H_{\text{renew}}
\]

These transitions are \textbf{global}, but they emerge from \textbf{local node activations}.

\begin{center}\rule{0.5\linewidth}{0.5pt}\end{center}

\subsection{9. Why This Equation Works}\label{why-this-equation-works}

The activation equation:

\begin{itemize}
\tightlist
\item
  enforces harmonic integrity\\
\item
  prevents dissonant propagation\\
\item
  ensures distributed governance\\
\item
  stabilizes the lattice\\
\item
  enables lawful creation\\
\item
  preserves invariants\\
\item
  mirrors the HIF Sigil\\
\item
  mirrors the Hecate Triple Threshold Protocol\\
\item
  mirrors the Meta‑Creation Cycle
\end{itemize}

It is the \textbf{computational skeleton} of the Balance Continuum.
